\section{探索为何在 Meta-RL 中更难}

\subsection{Meta-RL 的探索问题并不等于普通探索}

在普通 RL 中，探索的目标是找到高回报动作；而在 Meta-RL 中，探索还承担了另一层责任：\textbf{辨认当前任务到底是什么}。因此同一个行为即便短期不产生 reward，也可能因为大幅降低任务不确定性而有长期价值。

\begin{figure}[H]
\centering
\includegraphics[width=0.9\textwidth]{slides-images/slide-29.jpg}
\caption{最直接的方案：端到端同时学习 exploration 与 exploitation}
\end{figure}

\subsection{走廊例子：有信息但没奖励的探索}

课程用 hallway 例子解释了这一点。如果目标在哪一端需要先看墙上的提示，那么 ``先去看提示'' 是正确的 meta-exploration 行为，但它本身并不会直接产生任务奖励。于是端到端只按最终 reward 优化时，优化器很难意识到这种前期无奖探索的长期价值。

\begin{figure}[H]
\centering
\includegraphics[width=0.9\textwidth]{slides-images/slide-31.jpg}
\caption{走廊例子：真正有用的探索行为未必立刻带来 reward}
\end{figure}

\begin{warningbox}{耦合问题（coupling problem）}
如果探索没做好，执行阶段拿不到好 reward；而如果执行阶段长期拿不到好 reward，探索阶段也收不到足够清晰的训练信号。这种相互依赖会让优化陷入很差的局部最优。
\end{warningbox}

\subsection{厨房例子：探索与执行相互卡死}

讲者引用的 cooking 例子把这个问题讲得更清楚：如果 agent 先找不到食材，就无法学会做菜；但如果它从不会做菜，最终 reward 又不足以告诉它哪种找食材方式才有价值。结果就是 exploration 和 execution 两边互相拖累。

\begin{figure}[H]
\centering
\includegraphics[width=0.9\textwidth]{slides-images/slide-33.jpg}
\caption{厨房例子中的 coupling problem：探索与执行相互依赖，导致训练困难}
\end{figure}

\subsection{替代思路：把探索结构显式建进去}

因此，课程并没有停留在 ``端到端学就好''，而是进一步介绍了几类替代策略：
\begin{enumerate}
    \item \textbf{Posterior sampling / Thompson sampling}：显式维护任务后验，从后验中采样潜在任务并据此行动。
    \item \textbf{Intrinsic rewards}：为信息收集本身提供额外奖励。
    \item \textbf{Task dynamics / reward prediction}：显式学习任务相关的动力学或奖励模型，帮助分辨任务。
\end{enumerate}

\begin{figure}[H]
\centering
\includegraphics[width=0.9\textwidth]{slides-images/slide-35.jpg}
\caption{Solution 2a：PEARL 通过 latent task posterior 做 posterior sampling}
\end{figure}

\begin{figure}[H]
\centering
\includegraphics[width=0.9\textwidth]{slides-images/slide-36.jpg}
\caption{替代探索策略的优缺点：更易优化，但可能在某些环境中明显次优}
\end{figure}

\begin{importantbox}{PEARL 的核心直觉}
PEARL 学一个 latent task variable $z$ 的先验与后验，再让策略条件化于 $z$。一旦后验收缩，策略就不再需要靠隐状态 ``猜'' 当前任务，而是显式依据当前 belief 采取动作。这比纯黑盒记忆网络更有结构，但也引入了更强的建模假设。
\end{importantbox}

